% Appendix E — the disclosure boundary
\section{What is disclosed and what is protected}
\label{app:disclosure}

The language in which kernels are written, and the trained proposer, are proprietary. We protect them
by describing guarantees rather than mechanism; we do not protect anything a reader needs to check a
claim. \Cref{tab:disclosure} draws the line. The boundary does not make an unreleased component
reproducible, and we do not call it so: claims that depend on the proposer can be checked by
independent execution against frozen artifacts, and we will support such an audit.

\begin{table}[h]
\caption{The disclosure boundary.}
\label{tab:disclosure}
\centering\sffamily\footnotesize\lining
\begin{tabularx}{\linewidth}{@{}LL@{}}
\toprule
Disclosed & Protected \\
\midrule
Input and output contracts; encoder and policy dimensions and parameter counts & the kernel language's syntax and production rules \\
The six guarantees every kernel satisfies (\cref{sec:model}) & the operation catalog and its map to decoder tokens \\
Losses, objectives, optimizer settings and data separation & the decoding mask and how it enforces the guarantees \\
The game, the outcome $q$ and the incumbent rule & the teacher bank, the catalog and the search that found them \\
Every emitted program, in full & weights, training data and the emitters' source \\
The challenge, the entrants and every receipt & --- \\
A catalog-membership test, without the catalog & --- \\
\bottomrule
\end{tabularx}
\end{table}

The repository's published prose follows the same rule: the product is described as a model that reads
data and writes a program, and how candidate programs are generated is not described. Emitted code
necessarily reveals the computation it implements, which is why every program is published.
